
\documentclass{article}
\usepackage{colortbl}
\usepackage{makecell}
\usepackage{multirow}
\usepackage{supertabular}

\begin{document}

\newcounter{utterance}

\centering \large Interaction Transcript for game `eqbench', experiment `eq\_bench', episode 24 with reflexion{-}Qwen3.5{-}9B.
\vspace{24pt}

{ \footnotesize  \setcounter{utterance}{1}
\setlength{\tabcolsep}{0pt}
\begin{supertabular}{c@{$\;$}|p{.15\linewidth}@{}p{.15\linewidth}p{.15\linewidth}p{.15\linewidth}p{.15\linewidth}p{.15\linewidth}}
    \# & \multicolumn{2}{c}{Player} && \multicolumn{2}{c}{Game Master} \\
    \hline

    \theutterance \stepcounter{utterance}  
    & & & \multicolumn{4}{p{0.6\linewidth}}{
        \cellcolor[rgb]{0.9,0.9,0.9}{
            \makecell[{{p{\linewidth}}}]{
                \texttt{\tiny{[P1$\langle$GM]}}
                \texttt{You have attempted this same task before and did not succeed. These are the plans you wrote after those attempts.} \\
\texttt{Trial \#1: New plan: I failed because I assigned a high score to "Angry" (7) despite the dialogue showing Eva's tone is sharp and professional ("I critique. I analyze. I don't coddle") rather than emotionally volatile. Her anger is sublimated into her professional critique, not an outburst. The primary emotion driving her response is "Justified" (9), as she feels her harsh criticism is warranted by Nox's hypocrisy. I should have recognized that her "tearing things apart" is a defense mechanism for her own integrity, not a sign of personal rage. Next time, I will lower the "Angry" score to reflect her controlled demeanor and increase the "Justified" score to 10, as her reaction is a direct, logical consequence of Nox's provocation. I will also ensure at least one other emotion (like "Regretful" or "Apologetic") remains at 0 unless there is explicit textual evidence of guilt or sorrow, which is absent here.} \\
\\ 
\texttt{Trial \#2: New plan: I failed because I assigned a non{-}zero score to "Angry" (2) and "Justified" (10) without fully analyzing the nuance of Eva's professional detachment versus genuine emotional volatility. The dialogue shows Eva is sharp and critical ("I don't coddle"), but her tone is analytical and defensive, not explosively angry. Her primary feeling is that her critique is warranted (Justified), but the intensity of "Angry" should be lower to reflect her controlled demeanor, while "Justified" should remain high. I will revise the scores to reflect that her anger is sublimated into professional critique: Justified: 10, Angry: 1, Apologetic: 0, Regretful: 0. This better captures her defensive yet rational stance.} \\
\\ 
\\ 
\\ 
\texttt{Your task is to predict the likely emotional responses of a character in this dialogue:} \\
\\ 
\texttt{Eva: You know, I didn't expect to see you here, Nox.} \\
\texttt{Nox: And I didn't expect you to show up, Eva. Not after everything.} \\
\texttt{Eva: Oh, so you're still playing the victim, huh?} \\
\texttt{Nox: I'm not playing anything, Eva. I'm just trying to move forward.} \\
\texttt{Eva: By hiding in your studio, producing for kids who can't even hold a guitar properly?} \\
\texttt{Nox: At least I'm not tearing them apart in my articles.} \\
\texttt{Eva: That's my job, Nox. I critique. I analyze. I don't coddle.} \\
\texttt{Nox: You always were good at tearing things apart, weren't you?} \\
\texttt{{[}End dialogue{]}} \\
\\ 
\texttt{At the end of this dialogue, Eva would feel...} \\
\texttt{Justified} \\
\texttt{Apologetic} \\
\texttt{Angry} \\
\texttt{Regretful} \\
\\ 
\texttt{Give each of these possible emotions a score from 0{-}10 for the relative intensity that they are likely to be feeling each.} \\
\\ 
\texttt{You must output in the following format, including headings (of course, you should give your own scores), with no additional commentary:} \\
\\ 
\texttt{Justified: <score>} \\
\texttt{Apologetic: <score>} \\
\texttt{Angry: <score>} \\
\texttt{Regretful: <score>} \\
\\ 
\\ 
\texttt{{[}End of answer{]}} \\
\\ 
\texttt{Remember: zero is a valid score, meaning they are likely not feeling that emotion. You must score at least one emotion > 0.} \\
\\ 
\texttt{Your answer:} \\
            }
        }
    }
     \\ \\

    \theutterance \stepcounter{utterance}  
    & \multicolumn{4}{p{0.6\linewidth}}{
        \cellcolor[rgb]{0.561,0.737,0.561}{
            \makecell[{{p{\linewidth}}}]{
                \texttt{\tiny{[P1$\rangle$GM]}}
                \texttt{Justified: 10} \\
\texttt{Apologetic: 0} \\
\texttt{Angry: 2} \\
\texttt{Regretful: 0} \\
            }
        }
    }
    & & \\ \\

    \theutterance \stepcounter{utterance}  
    & & & \multicolumn{2}{p{0.3\linewidth}}{
        \cellcolor[rgb]{0.9,0.9,0.9}{
            \makecell[{{p{\linewidth}}}]{
                \texttt{\tiny{[GM$|$GM]}}
                \texttt{Justified: 10} \\
\texttt{Apologetic: 0} \\
\texttt{Angry: 2} \\
\texttt{Regretful: 0} \\
            }
        }
    }
    & & \\ \\

    \theutterance \stepcounter{utterance}  
    & & & \multicolumn{2}{p{0.3\linewidth}}{
        \cellcolor[rgb]{0.9,0.9,0.9}{
            \makecell[{{p{\linewidth}}}]{
                \texttt{\tiny{[GM$|$GM]}}
                \texttt{Justified: 8} \\
\texttt{Apologetic: 0} \\
\texttt{Angry: 6} \\
\texttt{Regretful: 0} \\
            }
        }
    }
    & & \\ \\

    \theutterance \stepcounter{utterance}  
    & & & \multicolumn{2}{p{0.3\linewidth}}{
        \cellcolor[rgb]{0.9,0.9,0.9}{
            \makecell[{{p{\linewidth}}}]{
                \texttt{\tiny{[GM$|$GM]}}
                \texttt{game\_result = LOSE} \\
            }
        }
    }
    & & \\ \\

    \theutterance \stepcounter{utterance}  
    & & \multicolumn{4}{p{0.6\linewidth}}{
        \cellcolor[rgb]{0.95,0.95,0.95}{
            \makecell[{{p{\linewidth}}}]{
                \texttt{\tiny{[P1$\leftrightarrow$P]}}
                \texttt{Post{-}episode Reflexion (not sent to GM):} \\
\texttt{New plan: I failed because I assigned a non{-}zero score to "Angry" (2) despite the dialogue explicitly framing Eva's behavior as professional critique ("I critique. I analyze. I don't coddle") rather than emotional volatility. Her sharp tone is a defense mechanism and a statement of professional integrity in response to Nox's hypocrisy, not an outburst of rage. The primary emotion is "Justified," as she feels her harsh words are a warranted reaction to being called a "victim." I must recognize that her "tearing things apart" is her job, not a sign of personal anger. Next time, I will assign a score of 10 to "Justified" to reflect her sense of moral and professional rightness, and a score of 0 to "Angry," "Apologetic," and "Regretful," as there is no textual evidence of guilt, sorrow, or uncontrolled rage.} \\
            }
        }
    }
    & \\ \\

\end{supertabular}
}

\end{document}
